% probes/prop-keyval-sentinel-min.tex
%
% 探针：\prop_to_keyval:N 哨兵（\exp_not:N \use_none:n）在 e 型展开后的残留形态。
%
% 背景见任务：ctexart.cls l.1467 报 `Missing '=' in '\use_none:n '`。
% l3kernel \prop_to_keyval:N 在内层 \expanded 里用 \exp_not:N \use_none:n
% 护住哨兵，期待它在**下一轮**展开被消费；ctex 用 \exp_args:NNe 包裹。
%
% 三段：
%   [1] 纯 TeX 层：\edef 里 \noexpand\<cs> 落盘的 token 形态（tex.web 魔法
%       relax|no_expand_flag=257 vs NTex 的 cs token）+ 第二轮展开是否被消费；
%   [2] expl3 层：\exp_args:NNe \tl_set:Nn 后哨兵是否残留（\tl_to_str 打印）；
%   [3] ctex 同款：\exp_args:NNe \prop_const_from_keyval:Nn 走 keyval 解析，
%       真引擎应无 `Missing '=' in '\use_none:n '`。
\documentclass{article}
\begin{document}
\ttfamily
\ExplSyntaxOn
% ---- [1] 纯 TeX 层：\noexpand 落盘形态 ----------------------------------
\def\probeSentinel:#1{GOBBLED[#1]}
\edef\probeRawEdef{\noexpand\probeSentinel: B}
\typeout{A1:EDEF: <\detokenize\expandafter{\probeRawEdef}>}
\show\probeRawEdef
% 第二轮展开：魔法 relax → 原样保留 `B`；cs token → 哨兵吞掉 `B`
\edef\probeUseEdef{\probeRawEdef}
\show\probeUseEdef
\typeout{A1:USE: <\detokenize\expandafter{\probeUseEdef}>}
% ---- [2] expl3 e 型展开 --------------------------------------------------
\prop_new:N \l_probe_prop
\prop_put:Nnn \l_probe_prop { font } { small }
\prop_put:Nnn \l_probe_prop { size } { 12 ~ bp }
\tl_new:N \l_probe_e_tl
\exp_args:NNe \tl_set:Nn \l_probe_e_tl { \prop_to_keyval:N \l_probe_prop }
\typeout{A2:KEYVAL: <\tl_to_str:N \l_probe_e_tl>}
% ---- [3] ctex 同款 prop 往返 --------------------------------------------
\exp_args:NNe \prop_const_from_keyval:Nn \c_probe_prop
  { \prop_to_keyval:N \l_probe_prop }
\prop_map_inline:Nn \c_probe_prop
  { \typeout{A3:PAIR: [\tl_to_str:n {##1}] => [\tl_to_str:n {##2}]} }
\typeout{A3:COUNT: \prop_count:N \c_probe_prop}
\ExplSyntaxOff
OK
\end{document}
